\section*{Test calculation}

Every timing uses bulk hcp cobalt with two atoms per cell
($a=\SI{2.507}{\angstrom}$, $c=\SI{4.069}{\angstrom}$). Spin--orbit coupling
uses SIESTA's off-site form~\citep{Cuadrado2012}, which keeps the coupling between orbitals on
different atoms. The pseudopotential is the fully relativistic PseudoDojo
nc-fr-04 PBE standard set~\citep{vanSetten2018}, read in PSML format~\citep{Garcia2018}, with 3s and 3p semicore states in the valence. The
basis is triple-$\zeta$ with polarization (energy shift \SI{0.005}{Ry}). That
gives 50 orbitals per cell, so each k point has a $100\times100$ complex
Hamiltonian once both spin components are included. The real-space mesh
cutoff is \SI{1200}{Ry}. Exchange and correlation use PBE~\citep{Perdew1996} with SIESTA's
standard noncollinear treatment of the gradient terms. Occupations follow
Fermi--Dirac smearing at \SI{2000}{K}. Both atomic moments start at
$1.6\,\mu_\mathrm{B}$, tilted \ang{45} from the $c$ axis, and are free to
relax. The high cutoff and temperature come from the magnetic-anisotropy
study this input was built for. They make the mesh work and the smearing larger
than in a typical production run.

An SCF run counts as converged when no density-matrix element changes by more
than $10^{-7}$ and no Hamiltonian element by more than \SI{e-4}{eV} between
iterations. Mixing is Pulay's scheme~\citep{Pulay1980} with weight 0.05 and eight stored steps. A tighter
Hamiltonian tolerance of \SI{e-5}{eV} cannot be met on every mesh. On the
$6\times6\times3$ mesh, unmodified SIESTA stalls with the largest change fixed
at \SI{1.4e-5}{eV} while the energy is constant to \SI{1}{\micro eV}. REMAG's GPU
runs did not converge within 400 iterations there or on the $14\times14\times8$
mesh. Convergence here
is SCF convergence at a fixed k mesh, not convergence in the number of k points.

The Brillouin zone is sampled with Monkhorst--Pack meshes~\citep{Monkhorst1976} of $6\times6\times3$,
$10\times10\times6$, $14\times14\times8$ and $20\times20\times11$ points. With
spin--orbit coupling SIESTA does not use time-reversal symmetry to reduce
them, so these give \num{108}, \num{600}, \num{1568} and \num{4400} k points.
Every k point is diagonalized in every SCF iteration.

\section*{What was timed}

We timed five configurations on the same input files.
\begin{enumerate}
\item Unmodified SIESTA 5.4.2 on one CPU. This is the numerical reference.
\item Unmodified SIESTA 5.4.2 with MPI on 28 processes. With \texttt{Diag.ParallelOverK}, each process diagonalizes its own share of the k points. The real-space mesh work is also shared out.
\item REMAG on 28 CPU threads. The k-point diagonalizations are shared among the threads. The mesh work stays on one thread.
\item REMAG on the GPU with its block-Jacobi eigensolver. After a Cholesky reduction, this solver repeatedly diagonalizes pairs of $32\times32$ blocks until the off-diagonal part is negligible. The SCF cycle, the initial density on the mesh and the final forces all run on the GPU. Groups of 64 k points are solved together, each starting from the previous iteration's eigenvectors.
\item REMAG on the GPU with a batched dense eigensolver. This is the same as the previous configuration, except for the solver. Each group of 256 k points is reduced with batched Cholesky and triangular solves and then diagonalized with cuSOLVER's \texttt{cusolverDnXsyevBatched}~\citep{cuSOLVER}.
\end{enumerate}
Everything runs in double precision. Unmodified SIESTA diagonalizes each k
point twice per iteration with LAPACK's divide-and-conquer solver. The first
solve gives eigenvalues to place the Fermi level, and the second gives
eigenvectors. In the timer records this is \num{35200} solves for
\num{4400} k points over four iterations. REMAG on the CPU keeps the
eigenvectors from one solve in memory. The REMAG settings used here are all
optional, and REMAG's defaults were not timed.

Each timing is the wall-clock time of one fresh process, from launch to exit.
It covers reading the input, setup of the basis and mesh, the SCF cycle, the
final forces and stress, and the output files. Runs never overlapped, and
repetitions of the five configurations were interleaved. The machine had one
NVIDIA H100 80GB HBM3 (driver 595.91.07, CUDA 12.9). The container was given
28 CPUs' worth of time on a host with two Intel Xeon Platinum 8480+ sockets,
112 cores and 224 hardware threads, shared with other tenants. Processes and
threads were not pinned to cores, and GPU runs used one CPU thread. GPU clocks
were not locked. All CPU code was built with GCC 13.3 against OpenBLAS 0.3.26
using one BLAS thread. The MPI build used Open MPI 4.1.6 and ScaLAPACK 2.2.1.
SIESTA's OpenMP support was off in every build.
